\chapter{Complexity — Convention}
%%% generated by the blueprint pipeline from TCSlib.Complexity.TuringMachine.Build.Convention
%%%%%%%%%%%%%%%%%%%%%%%%%%%%%%%%%%%%%%%%%%%%%%%%%%%%%%%%%%%%%%%%%%%%%%%%%%%%%%%
\section{Overview}

This module fixes the calling convention of the machine-construction library for
multi-tape Turing machines over the Boolean alphabet: the canonical seam configuration
(a designated anchor state, input head at its initial position, each work tape holding one
word written from the origin with its head at the origin, and empty output) at which loop
rounds and wrapped subroutines meet, together with the fact that a machine's initial
configuration is such a seam. It also introduces the pure list and arithmetic functions
against which the primitive machine contracts are stated: splitting a word at its last
$\mathrm{true}$, least solutions of padding length equations, fixed-width binary
increment with overflow, and the splitting off of a leading unary token.

%%%%%%%%%%%%%%%%%%%%%%%%%%%%%%%%%%%%%%%%%%%%%%%%%%%%%%%%%%%%%%%%%%%%%%%%%%%%%%%
\section{Declarations}

\begin{definition}[Seam configuration from tape words]
\lean{Turing.Cfg.ofWords}
\label{Turing.Cfg.ofWords}
\leanok
\uses{Turing.Cfg, Turing.FinTM.bufferTape}
Let $k \ge 0$, let $x$ be a Boolean input word, let $q$ be a state, and let
$w = (w_0, \dots, w_{k-1})$ be a family of Boolean words indexed by the $k$ work tapes.
The seam configuration $\mathrm{seam}(q, w)$ of a $k$-work-tape machine over the
Boolean alphabet on input $x$ is the configuration whose control is in the (non-halted)
state $q$, whose input head sits on the first input cell (position $1$), whose work tape
$i$ holds the word $w_i$ in cells $0, \dots, \abs{w_i} - 1$ and is blank at every other
integer cell, whose work-tape heads are all at cell $0$, and whose output tape is empty.
\end{definition}

\begin{lemma}[Initial configuration is a blank seam]
\lean{Turing.initCfg_ofWords}
\label{Turing.initCfg_ofWords}
\leanok
\uses{Turing.Cfg.init, Turing.Cfg.ofWords, Turing.FinTM.bufferTape_nil, Turing.MultiTapeTM, Turing.MultiTapeTM.initCfg}
\difficulty{1}
Let $M$ be a deterministic Turing machine with $k$ work tapes over the Boolean alphabet,
with start state $q_0$, and let $x$ be a Boolean input word. Then the initial
configuration of $M$ on $x$ equals the seam configuration $\mathrm{seam}(q_0, w)$
with $w_i = \varepsilon$ the empty word for every work tape $i$; that is, the
configuration in state $q_0$ with the input head on the first input cell, every work
tape entirely blank with its head at cell $0$, and empty output.
\end{lemma}

\begin{lemma}[Work tapes of a seam configuration]
\lean{Turing.Cfg.ofWords_workTapes}
\label{Turing.Cfg.ofWords_workTapes}
\leanok
\uses{Turing.Cfg.ofWords, Turing.FinTM.bufferTape}
\difficulty{1}
Let $k \ge 0$ be a number of work tapes, $x$ a Boolean input word, $q$ a state,
$w = (w_0, \dots, w_{k-1})$ a family of Boolean words, and $i < k$ a work-tape index.
In the seam configuration $\mathrm{seam}(q, w)$ of a $k$-work-tape machine over the
Boolean alphabet on input $x$ (state $q$, input head on the first input cell, work tape
$j$ holding $w_j$ from cell $0$, all work heads at cell $0$, empty output), work tape $i$
is exactly the buffer tape of $w_i$: for every integer cell $z$, its content is the
$z$-th letter of $w_i$ (counting from $0$) when $0 \le z < \abs{w_i}$, and blank
otherwise.
\end{lemma}

\begin{definition}[Split at the last true]
\lean{Turing.splitAtLastTrue}
\label{Turing.splitAtLastTrue}
\leanok
For a Boolean word $v$, $\mathrm{cut}(v)$ is the prefix of $v$ strictly
before its last occurrence of $\mathrm{true}$: if $v = u \, \mathrm{true} \,
\mathrm{false}^m$ for some word $u$ and $m \ge 0$, the result is $u$, and if $v$
contains no $\mathrm{true}$ (in particular if $v$ is empty), the result is undefined
(none).
\end{definition}

\begin{definition}[Solution of the polynomial padding equation]
\lean{Turing.solveSplit}
\label{Turing.solveSplit}
\leanok
For natural numbers $C$, $e$ and $n$, $\mathrm{solve}_{C,e}(n)$ is the least
$i \in \{0, \dots, n\}$ satisfying $i + C \cdot (i+1)^e = n$, or undefined (none) when
no such $i$ exists.
\end{definition}

\begin{definition}[Fixed-width binary increment]
\lean{Turing.incFixed}
\label{Turing.incFixed}
\leanok
Fixed-width little-endian binary increment with overflow, defined recursively on Boolean
words (least significant bit first): the empty word overflows (result none); a word
$\mathrm{false} \cdot r$ is sent to $\mathrm{true} \cdot r$; and a word
$\mathrm{true} \cdot r$ is sent to $\mathrm{false} \cdot r'$ where $r'$ is the
increment of $r$, overflowing whenever the increment of $r$ overflows. Thus a word of
length $\ell$ is mapped to the binary successor of the same length $\ell$, and the
result is none exactly when the word consists entirely of $\mathrm{true}$.
\end{definition}

\begin{definition}[Solution of a general padding equation]
\lean{Turing.solveSplitWith}
\label{Turing.solveSplitWith}
\leanok
For a function $f : \mathbb{N} \to \mathbb{N}$ and a natural number $n$,
$\mathrm{solve}_f(n)$ is the least $i \in \{0, \dots, n\}$ satisfying
$i + f(i) = n$, or undefined (none) when no such $i$ exists.
\end{definition}

\begin{definition}[Split off a leading unary token]
\lean{Turing.unaryTokenSplit}
\label{Turing.unaryTokenSplit}
\leanok
For a Boolean word $v$, $\mathrm{tok}(v)$ is the pair (token, remainder)
obtained by splitting off the maximal leading block of $\mathrm{true}$'s together with
the $\mathrm{false}$ that terminates it: if $v = \mathrm{true}^m \, \mathrm{false} \, r$
with $m \ge 0$, the result is $(\mathrm{true}^m \, \mathrm{false}, \, r)$, and if
$v = \mathrm{true}^m$ contains no $\mathrm{false}$ (including $v$ empty), the result is
$(\mathrm{true}^m, \varepsilon)$ with $\varepsilon$ the empty word.
\end{definition}
